% Section 12.1, from lectures/L12/appendix/a_realtime.md.

\section{Real Time}
\label{c12:app:a}

What the term claims, where latency comes from, and what \code{PREEMPT_RT} changes about the driver
you have spent eleven chapters writing.

\Secref{c12:app:b} is how to measure any of it without lying.

\subsection{The definition, and the trap in it}
\label{c12:app:a1}

\textbf{A real-time system is one in which being late is being wrong.} Not fast: correct within a
stated deadline, every time.

The trap is that ``real time'' sounds like a performance claim and is not. Consider:

\begin{narrowtable}{@{}lll@{}}
   & \thead{Average} & \thead{Worst observed} \\
  \midrule
  System A & 50 us & 30 ms \\
  System B & 5 ms & 5.1 ms \\
\end{narrowtable}

\noindent \textbf{System B is real time for a 6 ms deadline and system A is not}, despite being a
hundred times slower on average. A is faster; B is predictable. If your requirement is a 6 ms
deadline, A fails it several times an hour and no amount of tuning its average will help.

Everything in this chapter is about the right-hand column. \textbf{An average is not evidence}, and
a measurement without a stated worst case, a stated load and a stated duration is not a measurement
of anything.

Two more distinctions worth having:
\begin{itemize}
  \item \textbf{Hard real time}: a missed deadline is a system failure. An airbag, a motor
    commutation loop.
  \item \textbf{Soft real time}: a missed deadline degrades quality. Audio, video.
\end{itemize}
Linux with \code{PREEMPT_RT} is good soft real time and is used for hard real time by people who
have measured their own system very carefully. It is not a hard real-time operating system in the
sense that a small RTOS with a provable worst case is.

\subsection{Where latency comes from}
\label{c12:app:a2}

Three sources, and only three. Every real-time problem you will meet is one of them.

\textbf{Interrupts disabled.} Between \code{local_irq_disable()} and the matching enable, no
interrupt is taken, including yours. \code{spin_lock_irqsave} does this for the length of its
critical section, which is why \secref{c7:app:a5} insisted that those be short.

\textbf{Preemption disabled.} The interrupt is taken but the scheduler cannot switch to your task
afterwards. Holding any spinlock disables preemption, so a long critical section anywhere in the
kernel delays every high-priority task on that CPU.

\textbf{Priority inversion.} Your high-priority task is runnable but waiting on a lock held by a
low-priority task, which is itself not running because a medium-priority task is using the CPU. The
high-priority task now waits behind the medium one, which is the inversion.

The first two are code somebody wrote. The third is structural and needs a mechanism, which is
\S\ref{c12:app:a3}.

\subsection{Priority inversion, and Mars}
\label{c12:app:a3}

The canonical example is worth knowing accurately, because it is usually told badly.

Mars Pathfinder landed in 1997 and began resetting itself. A high-priority bus-management task
needed a mutex held by a low-priority meteorological task; a medium-priority communications task was
runnable and preempted the low-priority one; the low-priority task therefore never finished and
never released the mutex; the high-priority task missed its deadline; a watchdog concluded the
system had failed and reset it.

The fix was to enable \textbf{priority inheritance} on that mutex, and it was uploaded to Mars.

Priority inheritance is the mechanism: while a low-priority task holds a lock that a high-priority
task is waiting for, it \textbf{temporarily inherits the higher priority}, so it cannot be preempted
by anything in between. It finishes, releases the lock, and drops back.

The reason this story matters to a driver author is that the bug got through testing: it had caused
one or two resets in months of pre-flight testing, and nobody could reproduce or explain them. It
needed three tasks, one lock and a specific interleaving. It is the point \secref{c7:app:b} makes
about ordering bugs, one level up: \textbf{a system that has never missed a deadline is not a system
that cannot.}

\subsection{The four preemption models}
\label{c12:app:a4}

One question distinguishes them: \textbf{when a high-priority task becomes runnable, at what points
can the kernel switch to it?}

\begin{booktable}{@{}l>{\hsize=1.1\hsize}L>{\hsize=0.9\hsize}L@{}}
  \thead{Model} & \thead{Switches at} & \thead{For} \\
  \midrule
  \code{PREEMPT_NONE} & Only on return to userspace, or a voluntary sleep & Throughput. Servers \\
  \code{PREEMPT_VOLUNTARY} & Those, plus explicit \code{might_sleep} points & Desktops, historically \\
  \code{PREEMPT} & Those, plus almost anywhere in the kernel & Desktops and embedded. \textbf{This course's default} \\
  \code{PREEMPT_RT} & Those, plus inside most spinlock sections & Real time \\
\end{booktable}

\noindent The progression is one of removing places where the kernel refuses to be interrupted. Each
step costs throughput, because each adds preemption points that must be checked, and buys a lower
worst case.

\textbf{\code{PREEMPT_RT} was merged into mainline in 6.12}, which is the kernel this course builds.
That is recent and it matters: before it, running a real-time Linux meant applying an out-of-tree
patch series that trailed the kernel it patched, sometimes did not exist for a given version, and
was a standing porting cost. \file{kernel/rt.config} is now \textbf{four lines}, and the whole of
the difference between the two kernels this chapter compares.

One of those four lines is \code{CONFIG_EXPERT=y}, and it is not decoration: \code{PREEMPT_RT} is
declared \code{depends on EXPERT && ARCH_SUPPORTS_RT}, so without it the symbol is never offered and
\code{CONFIG_PREEMPT_RT=y} is silently discarded, leaving a kernel that builds, boots, says
\code{PREEMPT} and is not real time. That is exactly the hidden-prompt mechanism
Exercise~\ref{c3:ex:3} makes you work through with \code{GPIO_SYSFS}.

\subsection{What \code{PREEMPT_RT} actually changes}
\label{c12:app:a5}

Four things, and the first is the one that reaches your code.

\textbf{Most spinlocks become sleeping locks.} \code{spinlock_t} is replaced by an \code{rt_mutex}
underneath, so code holding one can be preempted and can even sleep. That removes the largest source
of preemption-disabled time in the kernel at a stroke.

It also \textbf{changes what your driver is allowed to assume.} Code that took a \code{spinlock_t}
and relied on preemption being disabled, for instance to touch per-CPU data safely, is wrong under
\code{PREEMPT_RT} and was correct before.

\textbf{Most interrupt handlers become threads.} Every handler registered with \code{request_irq}
runs in a kernel thread by default, exactly like the threaded handlers \secref{c8:app:b4}
introduced, unless it asks not to with \code{IRQF_NO_THREAD}; per-CPU interrupts, and a primary
handler requested with \code{IRQF_ONESHOT}, stay in hard interrupt context as well. So the
hard-interrupt path shrinks to almost nothing and your handler becomes schedulable, preemptible and
prioritisable.

\textbf{Mutexes get priority inheritance}, which is \secref{c12:app:a3}'s fix applied by default.

\textbf{Softirqs are threaded too}, removing another unbounded source of delay.

\subsubsection{What does not change}
\label{c12:app:a:what-does-not-change}

\code{raw_spinlock_t} still spins, still disables preemption, and still cannot sleep. It exists for
the places that genuinely cannot be preempted: the scheduler itself, the timer core, the interrupt
entry path.

\textbf{Reaching for \code{raw_spinlock_t} to make your driver faster reintroduces exactly the
latency \code{PREEMPT_RT} removed}, and does so in a way that is invisible until somebody measures
the worst case. A driver should essentially never contain one.

\subsection{What makes a driver hostile to real time}
\label{c12:app:a6}

Four habits. Look for them in your own driver; at least two are in it.

\textbf{Long critical sections.} Anything held across a loop, a copy, or a call into unfamiliar
code. The fix is to shorten it, not to change the lock type.

\textbf{Interrupts disabled across real work.} \code{spin_lock_irqsave} around a page of processing
rather than around the three lines that touch shared state.

\textbf{Busy-waiting.} \code{udelay(50)} is legal, works, and spends fifty microseconds of a CPU
doing nothing. \Secref{c9:app:b3} is where the alternatives are; the rule is that \code{udelay} is
for a few microseconds of genuine hardware settling time and nothing else.

\textbf{Per-CPU assumptions.} Code that updates a per-CPU variable while holding a
\code{spinlock_t}, relying on preemption being disabled to keep every other task on that CPU away
from it. Correct on \code{PREEMPT}, wrong on \code{PREEMPT_RT}, where the lock still keeps the task
on its CPU but no longer keeps other tasks off it. \code{local_lock_t} is the lock that does that on
both kernels.

\subsection{Priorities, and the throttle that saves you}
\label{c12:app:a7}

Making a task real-time is one call, and the consequences deserve a paragraph.

\begin{booktable}{@{}lL@{}}
  \thead{Policy} & \thead{Behaviour} \\
  \midrule
  \code{SCHED_OTHER} & The normal fair scheduler. Everything by default \\
  \code{SCHED_FIFO} & Runs until it blocks or a higher priority becomes runnable \\
  \code{SCHED_RR} & The same, with a time slice between equal priorities \\
  \code{SCHED_DEADLINE} & Declares a period and a budget; the kernel admits or refuses \\
\end{booktable}

\noindent \textbf{A \code{SCHED_FIFO} task that does not block will not be preempted by anything of
lower priority, including the shell you would use to kill it.} On a single-CPU machine that is an
unrecoverable hang, and it is the classic first mistake.

Linux therefore throttles: \code{sched_rt_runtime_us} defaults to 950,000 in every 1,000,000, so
real-time tasks get at most 95\% of a CPU and something else can always run. That saves you, and it
also means \textbf{a real-time task can be preempted by the throttle}, which people discover as an
unexplained 50 ms gap. Turning the throttle off is possible and should be a deliberate, documented
decision.

\code{SCHED_DEADLINE} is the interesting one for periodic work: you state a period and a worst-case
execution time and the kernel performs an admission test, refusing the request if the deadline
tasks together would need more CPU time than the system allows them. On one CPU that is a real
guarantee rather than a priority; on several it bounds how late a task can be rather than promising
it never is, which is still more than a priority says.

\textbf{Two tools that belong next to this, and that this course does not teach.} On a real system
you would also reach for \textbf{CPU isolation} (\code{isolcpus}, \code{nohz_full}) to keep
everything else off the core your task runs on, and \textbf{IRQ affinity}
(\file{/proc/irq/N/smp_affinity}) to decide which core takes the interrupt. Both are absent here,
and the reason is the target rather than the topic: this course runs on a two-CPU QEMU guest, and
isolating one of two emulated CPUs from a scheduler that is itself emulated produces a number that
says nothing about a real machine. That is the same standard \secref{c12:app:b6} applies to
everything else measured here. Know that the two knobs exist and that they are usually the next
thing tried after \code{PREEMPT_RT}; measure them somewhere with real cores.

\subsection{The comparison, measured}
\label{c12:app:a8}

The lab measures handler-to-userspace latency, the interval \secref{c12:app:b3} can give as an
absolute figure, on both kernels, idle and under four spinning processes on two CPUs. One run of
each:

\begin{narrowtable}{@{}llrr@{}}
  \thead{Kernel} & \thead{Load} & \thead{Mean} & \thead{Worst} \\
  \midrule
  \code{PREEMPT} & idle & 449 us & \textbf{11,321 us} \\
  \code{PREEMPT_RT} & idle & 551 us & \textbf{4,964 us} \\
  \code{PREEMPT} & busy & 1,685 us & \textbf{29,729 us} \\
  \code{PREEMPT_RT} & busy & 2,393 us & \textbf{18,559 us} \\
\end{narrowtable}

\noindent Read the two columns separately, because they disagree.

\textbf{The mean is worse under \code{PREEMPT_RT}}, by about 23\% idle and 42\% loaded. That is the
cost of the extra preemption points, the threaded handlers and the inheritance bookkeeping, and it
is expected.

\textbf{The worst case is better under \code{PREEMPT_RT}}, by a factor of 2.3 idle and 1.6 loaded.

\textbf{That is the whole trade, and it is the definition of real time in one table.} A change that
made the average worse and the maximum better would be a straightforward regression on a throughput
system and is exactly what you are buying here.

\textbf{And every number above is worthless on its own}, for reasons \secref{c12:app:b} sets out in
full. They were measured under an emulator, over 2,000 samples, on one run. Repeating the idle
\code{PREEMPT} measurement gave a worst case of 2,502 us rather than 11,321. The \emph{shape} of the
comparison survives that; the values do not.

\lecturefigure[0.96]{lectures/L12/appendix/images/latency.png}
  {Two panels of grouped bars. On the left, mean latency, where PREEMPT is lower at 449 against 551
   microseconds idle and 1,685 against 2,393 busy. On the right, the worst case, where PREEMPT\_RT
   is lower at 4,964 against 11,321 idle and 18,559 against 29,729 busy, annotated with the same
   PREEMPT measurement repeated at 2,502.}
  {fig:latency}
